\chapter{Exposition de contrepartie}\label{ch:rc:counterparty-exposure}

Lorsque Lehman Brothers a déposé son bilan le 15 septembre 2008, elle était partie à environ
930\,000 transactions dérivées, selon son propre premier décompte. Chaque contrepartie détenait
désormais des contrats avec une entreprise en défaut~: elle pouvait les résilier, compenser ce
qu'elle devait par ce qui lui était dû, conserver le collatéral qu'elle détenait et remplacer les couvertures sur un marché où toutes
les autres contreparties entraient du même côté au même moment. En novembre, la plupart
des contrats avaient été résiliés~; vingt mois plus tard, les créances des banques au titre des pertes sur dérivés
dépassaient 50 milliards de dollars. Ce qu'une banque risque de perdre lorsqu'une contrepartie fait défaut
dépend de la valeur que prendront les contrats à ce moment-là, valeur inconnue aujourd'hui~: la perte est une
option sur la valeur future d'un portefeuille. Ce chapitre la mesure. Il ouvre la partie~III, dont les
chapitres transforment l'exposition en prix (ajustements de valeur pour risque de crédit et pour risque de crédit propre), en
coûts de financement, de marge et de fonds propres, et en un desk qui les gère.

\section{Exposition et profils}

\begin{definition}[Risque de crédit de contrepartie]\label{def:rc:counterparty-exposure:ccr}
Le \emph{risque de crédit de contrepartie}\index{risque de crédit de contrepartie} est le risque que l'autre
partie à un dérivé fasse défaut avant le paiement final du contrat alors que celui-ci a une
valeur positive pour la partie survivante. Contrairement au risque de crédit d'un prêt, le montant
exposé est incertain et peut changer de signe.
\end{definition}

\begin{definition}[Exposition de contrepartie]\label{def:rc:counterparty-exposure:exposure}
L'\emph{exposition de contrepartie}\index{exposition de contrepartie} à une date future $t$ est la
perte si la contrepartie fait défaut à ce moment et que rien n'est recouvré~: $E_t = \max(V_t-C_t,0)$,
où $V_t$ est la valeur des contrats pour la banque et $C_t$ le collatéral qu'elle peut conserver.
\end{definition}

\begin{definition}[Profils d'exposition]\label{def:rc:counterparty-exposure:profiles}
L'\emph{exposition attendue}\index{exposition attendue} est $\mathrm{EE}(t) = \E[E_t]$~;
l'\emph{exposition négative attendue}\index{exposition négative attendue} est
$\mathrm{ENE}(t) = \E[\min(V_t-C_t,0)]$, l'exposition attendue de la contrepartie envers la banque~;
l'\emph{exposition future potentielle}\index{exposition future potentielle} $\mathrm{PFE}_q(t)$ est le
quantile d'ordre $q$ de $E_t$ (97.5~\% ici)~; l'\emph{exposition positive
attendue}\index{exposition positive attendue} sur un horizon $H$ est la moyenne temporelle
$\mathrm{EPE} = \frac1H\int_0^H\mathrm{EE}(t)\,dt$.
\end{definition}

Les EE et les ENE permettent de chiffrer les ajustements des chapitres suivants~; la PFE fixe les limites de crédit (le
montant maximal que la banque pourrait perdre avec un niveau de confiance donné)~; l'EPE est la mesure
synthétique unique qui sous-tend les mesures réglementaires d'exposition.

\begin{example}[Deux transactions avec une même contrepartie]\label{ex:rc:counterparty-exposure:profiles}
Une banque reçoit le taux fixe au pair de 3.80~\% sur un swap décennal de 100~millions de USD
(courbe SOFR du chapitre~1, Hull--White avec $\kappa = 3~\%$ et une volatilité normale de 90 points de
base), et avec la même contrepartie, reçoit le taux fixe en USD à 3.80~\% sur 110~millions de USD, et
paie le taux fixe en EUR au taux au pair ESTR de 2.59~\% sur 100~millions d'EUR pendant dix ans, en
échangeant les notionnels à la fin (EURUSD à 1.10, volatilité de 8~\%). Sur 4\,000 trajectoires simulées
tous les dix jours ouvrés, l'EE du swap culmine à 3.13~millions de USD après trois ans et sa PFE
à 18.26~millions de USD après 3.9 ans~: le risque de taux croît avec le temps tandis que la duration
restante diminue. L'EE du swap de devises monte à 7.09~millions de USD et sa PFE à
39.93~millions de USD juste avant l'échéance, sous l'effet de l'échange final
(\cref{fig:rc:counterparty-exposure:profiles}).
\end{example}

\begin{omfigure}
\begin{tikzpicture}
\begin{axis}[width=11.5cm,height=5.4cm,
  xlabel={temps (années)},ylabel={millions de USD},
  tick label style={font=\small},label style={font=\small},
  xmin=0,xmax=10,ymin=0,ymax=42,grid=major,grid style={gray!25},
  legend columns=2,legend cell align=left,legend style={font=\footnotesize,at={(0.5,-0.25)},anchor=north,draw=none,fill=none,/tikz/every even column/.append style={column sep=8pt}},
  table/col sep=comma]
\addplot[omDef,very thick] table[x=t,y=swap_ee]{figdata/rates-credit-risk/17-counterparty-exposure/profiles.csv};
\addplot[omDef,thick,dashed] table[x=t,y=swap_pfe]{figdata/rates-credit-risk/17-counterparty-exposure/profiles.csv};
\addplot[omThm,very thick] table[x=t,y=xccy_ee]{figdata/rates-credit-risk/17-counterparty-exposure/profiles.csv};
\addplot[omThm,thick,dashed] table[x=t,y=xccy_pfe]{figdata/rates-credit-risk/17-counterparty-exposure/profiles.csv};
\legend{EE du swap,PFE du swap (97.5~\%),EE du swap de devises,PFE du swap de devises (97.5~\%)}
\end{axis}
\end{tikzpicture}

\omcaption{Profils d'exposition d'un swap de taux décennal (en bosse~: l'incertitude croît
tandis que les flux restants s'éteignent) et d'un swap de devises décennal (croissant jusqu'à
l'échéance~: l'échange final des notionnels maintient le risque de change jusqu'au bout). Les petites dents
correspondent aux paiements annuels. Données~: le tutoriel du chapitre.}\label{fig:rc:counterparty-exposure:profiles}
\end{omfigure}

\section{Compensation et collatéral}

\begin{definition}[Ensemble de compensation, compensation avec déchéance du terme]\label{def:rc:counterparty-exposure:netting}
Un \emph{ensemble de compensation}\index{ensemble de compensation} est le groupe des transactions conclues
avec une même contrepartie au titre d'une même convention-cadre juridiquement opposable. La \emph{compensation avec déchéance du terme}\index{compensation avec déchéance du terme} (close-out netting)
permet à la partie survivante de résilier chaque transaction de l'ensemble de compensation en cas de défaut et de
compenser les valeurs entre elles en un montant net unique, de sorte que l'exposition vaut
$\max(\sum_iV_i,0)$, et non $\sum_i\max(V_i,0)$.
\end{definition}

\begin{dated}{2026-09}{Avis juridiques sur la compensation}\label{dat:rc:counterparty-exposure:opinions}
Le fait que la \omterm{def:rc:counterparty-exposure:netting}{compensation avec déchéance du terme} soit valable relève du droit de l'insolvabilité de chaque contrepartie. L'ISDA a
publié des avis juridiques sur l'applicabilité de la compensation au titre de ses conventions-cadres couvrant plus de 90
juridictions, et sur le collatéral couvrant plus de 60, mis à jour en général chaque année.
\end{dated}

\begin{example}[Le bénéfice de la compensation]\label{ex:rc:counterparty-exposure:netting}
Une fois compensées, les deux transactions de l'\cref{ex:rc:counterparty-exposure:profiles} ont une EE qui culmine
à 7.77~millions de USD après 3.9 ans, alors que les deux EE isolées s'additionnent à 9.24~millions de USD~:
une économie de 16~\% (\cref{fig:rc:counterparty-exposure:netting}). L'ENE de l'\omterm{def:rc:counterparty-exposure:netting}{ensemble de compensation} atteint
$-\text{USD}~21.0$ millions~: la banque paie dans la devise à taux le plus bas, et l'échange final
se fait à un taux à terme supérieur au taux au comptant, de sorte que la transaction dérive en faveur de la contrepartie.
\end{example}

\begin{omfigure}
\begin{tikzpicture}
\begin{axis}[width=11.5cm,height=5.4cm,
  xlabel={temps (années)},ylabel={millions de USD},
  tick label style={font=\small},label style={font=\small},
  xmin=0,xmax=10,ymin=-25,ymax=12,grid=major,grid style={gray!25},
  legend columns=3,legend style={font=\footnotesize,at={(0.5,-0.25)},anchor=north,draw=none,fill=none,/tikz/every even column/.append style={column sep=8pt}},
  table/col sep=comma]
\addplot[omExo,thick,dashed] table[x=t,y=sum_ee]{figdata/rates-credit-risk/17-counterparty-exposure/netting.csv};
\addplot[omDef,very thick] table[x=t,y=net_ee]{figdata/rates-credit-risk/17-counterparty-exposure/netting.csv};
\addplot[omThm,very thick] table[x=t,y=net_ene]{figdata/rates-credit-risk/17-counterparty-exposure/netting.csv};
\legend{somme des EE isolées,EE compensée,ENE compensée}
\end{axis}
\end{tikzpicture}

\omcaption{Compensation du swap et du swap de devises~: l'EE compensée se situe sous la somme des
EE isolées, et l'ENE compensée est plus grande en valeur absolue que l'EE parce que l'\omterm{def:rc:counterparty-exposure:netting}{ensemble de compensation}
dérive en faveur de la contrepartie. Données~: le tutoriel du chapitre.}\label{fig:rc:counterparty-exposure:netting}
\end{omfigure}

La credit support annex (CSA) du livre~2, chapitre~10 oblige les parties à échanger
de la marge de variation à mesure que la valeur de l'\omterm{def:rc:counterparty-exposure:netting}{ensemble de compensation} évolue. Trois termes en fixent le montant.

\begin{definition}[Seuil de collatéralisation, montant minimal de transfert, montant indépendant]\label{def:rc:counterparty-exposure:csa}
Le \emph{seuil de collatéralisation}\index{seuil de collatéralisation} (collateral threshold) $\mathrm{Th}$ est le niveau d'exposition
qu'une partie accepte avant d'appeler du collatéral~: le collatéral appelé vaut $V-\mathrm{Th}$ au-delà de ce seuil. Le
\emph{montant minimal de transfert}\index{montant minimal de transfert} (minimum transfer amount) (MTA) est la plus petite variation
de collatéral qui fait l'objet d'un appel. Le \emph{montant indépendant}\index{montant indépendant} (independent amount) est
un collatéral déposé indépendamment de la valeur, une réserve contre les mouvements pendant le temps nécessaire
pour clôturer les positions.
\end{definition}

\section{Simulation de l'exposition}

\begin{method}[Exposition par simulation]\label{met:rc:counterparty-exposure:sim}
(1) Choisir une grille temporelle incluant les dates de paiement des transactions. (2) Simuler les facteurs
de risque conjointement sur la grille (ici le facteur de Hull--White simulé exactement, et EURUSD autour de son
forward). (3) Réévaluer chaque transaction sur chaque trajectoire et à chaque date, en forme fermée lorsque c'est possible, par
régression sinon (la méthode de Longstaff--Schwartz du livre~5). (4) Sommer chaque
\omterm{def:rc:counterparty-exposure:netting}{ensemble de compensation}~; appliquer l'accord de collatéralisation trajectoire par trajectoire. (5) Calculer la moyenne et les quantiles
sur les trajectoires.
\end{method}

\omcode{code/firm/exposure/firm_exposure.py}{58}{78}{Génération des scénarios~: le facteur de Hull--White sur la grille, avec sa dérive et sa variance exactes, et EURUSD autour de son forward.}\label{lst:rc:counterparty-exposure:sim}

Le coût est le produit du nombre de trajectoires, de dates, de transactions et du coût d'une valorisation~: des millions de
valorisations pour un seul \omterm{def:rc:counterparty-exposure:netting}{ensemble de compensation}, des milliards pour une banque. L'étape de réévaluation est celle où
les moteurs passent l'essentiel de leur temps~; l'étape d'agrégation (compensation, collatéral, profils) porte sur
de très grands tableaux et c'est la partie que la banque écrit en C++20 et en Rust.

\section{La période de marge en risque et le profil collatéralisé}

Le collatéral ne supprime pas l'exposition~: lorsqu'une contrepartie fait défaut, le dernier collatéral
qu'elle a déposé reflète la valeur lors du dernier appel de marge abouti, et la partie survivante
a besoin de temps pour constater le défaut, contester, résilier et reconstituer sa couverture.

\begin{definition}[Période de marge en risque]\label{def:rc:counterparty-exposure:mpor}
L'exposition au titre d'un accord de collatéralisation à $t$ est $\max(V_t-C_{t-\delta},0)$, où $\delta$ est la
période de marge en risque du livre~1~: le collatéral est celui convenu $\delta$ auparavant.
L'exposition est la variation de la valeur de l'\omterm{def:rc:counterparty-exposure:netting}{ensemble de compensation} sur $\delta$, augmentée du seuil et du
\omterm{def:rc:counterparty-exposure:csa}{montant minimal de transfert}.
\end{definition}

\omcode{code/firm/exposure/firm_exposure.py}{149}{165}{Le collatéral détenu au titre d'une CSA bilatérale, et l'exposition qui utilise le collatéral convenu une période de marge en risque plus tôt.}\label{lst:rc:counterparty-exposure:csa}

\begin{omfigure}
\begin{tikzpicture}[x=1.05cm,y=1cm,font=\small]
\draw[->,thick] (0,0) -- (10.8,0) node[right] {$t$};
\foreach \x in {1,2,3,4} {\draw[omDef,thick] (\x,0.12) -- (\x,-0.12);}
\node[omDef,below,align=center] at (2.5,-0.2) {appels de marge quotidiens honorés};
\draw[omThm,very thick] (5,0.3) -- (5,-0.3);
\node[omThm,above,align=center] at (5,0.35) {dernier appel honoré\\$C = V_{t-\delta}$};
\draw[omThm,very thick,dashed] (7,0.3) -- (7,-0.3);
\node[omThm,above] at (7,0.35) {défaut};
\draw[omThm,very thick] (9.6,0.3) -- (9.6,-0.3);
\node[omThm,above,align=center] at (9.6,0.35) {résiliation\\$V_t$};
\draw[<->,omExo,thick] (5,-0.75) -- (9.6,-0.75);
\node[omExo,below,align=center] at (7.3,-0.8) {période de marge en risque $\delta$~: aucun mouvement de collatéral~;\\exposition $= \max(V_t - C,\,0)$};
\end{tikzpicture}

\omcaption{La période de marge en risque~: le dernier collatéral reçu correspond à la valeur
au dernier appel réussi~; entre celui-ci et la résiliation, la contrepartie ne dépose rien, et
la valeur continue de varier. Schéma de principe.}\label{fig:rc:counterparty-exposure:mpor}
\end{omfigure}

\begin{example}[Collatéralisation de l'ensemble de compensation]\label{ex:rc:counterparty-exposure:csa}
Avec une CSA bilatérale à seuil nul et un \omterm{def:rc:counterparty-exposure:csa}{montant minimal de transfert} de 500\,000~dollars,
l'EE de l'\omterm{def:rc:counterparty-exposure:netting}{ensemble de compensation} tombe à un niveau presque plat de 1.30~million de dollars à son maximum, et sa PFE à
6.45~millions de dollars avec une période de marge en risque de dix jours ouvrés~; l'EPE sur la durée de vie
passe de 6.44~millions de dollars à 0.95~million de dollars, soit 15~\% du chiffre sans collatéral. Avec
vingt jours, l'EE maximale est de 1.83~million de dollars, soit $\sqrt2$ fois plus élevée~; avec un seuil de
10~millions de dollars, elle est de 3.81~millions de dollars (la \cref{fig:rc:counterparty-exposure:collat}).
\end{example}

\begin{omfigure}
\begin{tikzpicture}
\begin{axis}[width=11.5cm,height=5.4cm,
  xlabel={temps (années)},ylabel={exposition attendue (millions de dollars)},
  tick label style={font=\small},label style={font=\small},
  xmin=0,xmax=10,ymin=0,ymax=9,grid=major,grid style={gray!25},
  legend columns=2,legend cell align=left,legend style={font=\footnotesize,at={(0.5,-0.25)},anchor=north,draw=none,fill=none,/tikz/every even column/.append style={column sep=8pt}},
  table/col sep=comma]
\addplot[omExo,very thick] table[x=t,y=none]{figdata/rates-credit-risk/17-counterparty-exposure/collat.csv};
\addplot[omMeth,very thick] table[x=t,y=th10]{figdata/rates-credit-risk/17-counterparty-exposure/collat.csv};
\addplot[omThm,thick] table[x=t,y=csa20]{figdata/rates-credit-risk/17-counterparty-exposure/collat.csv};
\addplot[omDef,thick] table[x=t,y=csa10]{figdata/rates-credit-risk/17-counterparty-exposure/collat.csv};
\legend{sans CSA,seuil de 10~millions de dollars (10 jours),seuil nul (20 jours),seuil nul (10 jours)}
\end{axis}
\end{tikzpicture}

\omcaption{\omterm{def:rc:counterparty-exposure:profiles}{Exposition attendue} de l'\omterm{def:rc:counterparty-exposure:netting}{ensemble de compensation} sans collatéral et sous trois CSA. Une
CSA à seuil nul ramène le profil aux variations de valeur survenant pendant la période de marge en risque~;
doubler cette période les multiplie par environ $\sqrt2$~; un seuil ajoute de l'exposition à hauteur de
sa taille. Les crans à chaque anniversaire correspondent aux paiements annuels, que le collatéral
retardé ne reflète pas encore. Données~: le tutoriel du chapitre.}\label{fig:rc:counterparty-exposure:collat}
\end{omfigure}

\begin{dated}{2026-09}{Planchers de la période de marge en risque}\label{dat:rc:counterparty-exposure:mpor}
L'approche standard pour le \omterm{def:rc:counterparty-exposure:ccr}{risque de crédit de contrepartie} (SA-CCR) du Comité de Bâle (mars 2014)
fixe une période de marge en risque minimale de dix jours ouvrés pour les dérivés bilatéraux avec
marge quotidienne, cinq pour les transactions clients compensées et vingt pour les \omterm{def:rc:counterparty-exposure:netting}{ensembles de compensation} de 5\,000 transactions ou plus
non compensées par une contrepartie centrale, doublée pour les \omterm{def:rc:counterparty-exposure:netting}{ensembles de compensation} présentant des
litiges en cours~; avec une marge appelée tous les $N$ jours, elle est de $10+N-1$ jours.
\end{dated}

\section{Risque de corrélation défavorable}

\begin{definition}[Risque de corrélation défavorable]\label{def:rc:counterparty-exposure:wwr}
Le \emph{risque de corrélation défavorable}\index{risque de corrélation défavorable} (wrong-way risk) est une dépendance positive entre la
\omterm{def:rc:reduced-form-credit:pd}{probabilité de défaut} d'une contrepartie et l'exposition de la banque à celle-ci~: l'exposition est
la plus élevée lorsque le défaut est le plus probable. Il est \emph{spécifique} lorsque la dépendance est intégrée à
la transaction (protection achetée auprès d'une contrepartie sur elle-même ou sur son propre pays) et
\emph{général} lorsqu'il passe par des facteurs de marché communs.
\end{definition}

\begin{example}[Une contrepartie qui s'affaiblit avec l'euro]\label{ex:rc:counterparty-exposure:wwr}
Soit la contrepartie du swap de devises d'intensité de défaut $2~\%\times(X_t/F_X(0,t))^{-5}$~:
une baisse de 10~\% de l'euro en dessous de son cours à terme fait monter son intensité d'environ 70~\%. L'exposition de la
banque est la plus élevée lorsque l'euro a baissé. En pondérant les trajectoires par la \omterm{def:rc:reduced-form-credit:pd}{probabilité de
défaut} de chaque période, l'\omterm{def:rc:counterparty-exposure:profiles}{exposition attendue} conditionnelle au défaut vaut 2.52 fois l'EE
inconditionnelle à son maximum, et 2.22 fois en moyenne sur la durée de vie
(la \cref{fig:rc:counterparty-exposure:wwr}).
\end{example}

\begin{omfigure}
\begin{tikzpicture}
\begin{axis}[width=11.5cm,height=5.4cm,
  xlabel={temps (années)},ylabel={millions de USD},
  tick label style={font=\small},label style={font=\small},
  xmin=0,xmax=10,ymin=0,ymax=20,grid=major,grid style={gray!25},
  legend columns=2,legend cell align=left,legend style={font=\footnotesize,at={(0.5,-0.25)},anchor=north,draw=none,fill=none,/tikz/every even column/.append style={column sep=8pt}},
  table/col sep=comma]
\addplot[omDef,very thick] table[x=t,y=ee]{figdata/rates-credit-risk/17-counterparty-exposure/wwr.csv};
\addplot[omThm,very thick,dashed] table[x=t,y=cee]{figdata/rates-credit-risk/17-counterparty-exposure/wwr.csv};
\legend{EE,EE sachant le défaut (wrong-way)}
\end{axis}
\end{tikzpicture}

\omcaption{\omterm{def:rc:counterparty-exposure:profiles}{Exposition attendue} du swap de devises, inconditionnellement et conditionnellement au
défaut de la contrepartie à chaque date, lorsque l'intensité de la contrepartie augmente quand l'euro
baisse. Les expositions qui comptent sont celles des scénarios où la contrepartie
fait défaut. Données~: le tutoriel du chapitre.}\label{fig:rc:counterparty-exposure:wwr}
\end{omfigure}

\begin{dated}{2026-09}{Les dérivés de Lehman}\label{dat:rc:counterparty-exposure:lehman}
Selon la Financial Crisis Inquiry Commission, Lehman avait d'abord revendiqué environ
930\,000 transactions dérivées à sa faillite~; la plupart avaient été résiliées au plus tard le
13 novembre 2008, 18\,000 étaient encore en cours en janvier 2009, et en mai 2010 les banques avaient
déposé plus de 50 milliards de dollars de créances au titre des pertes sur dérivés envers Lehman.
\end{dated}

\section{Tutoriel~: exposition d'un ensemble de compensation}\label{tut:rc:counterparty-exposure:tutorial}

\begin{tutorial}
\textbf{Objectif.} Simuler l'exposition d'un swap, d'un swap de devises et de leur \omterm{def:rc:counterparty-exposure:netting}{ensemble de compensation},
avec et sans marge de variation. \textbf{État final~:} les résultats chiffrés des
\cref{ex:rc:counterparty-exposure:profiles,ex:rc:counterparty-exposure:netting,ex:rc:counterparty-exposure:csa,ex:rc:counterparty-exposure:wwr}
et leurs quatre graphiques.
\begin{tutsteps}
\item \textbf{Scénarios}~: \texttt{simulate(MKT, 10, 26, 4000)}~; vérifier qu'EURUSD est une
martingale autour de son cours à terme.
\item \textbf{Valeurs}~: \texttt{Swap.values(sc)}, \texttt{XccySwap.values(sc)}~; les deux
valent zéro au départ.
\item \textbf{Profils}~: \texttt{aggregate([swap])}, l'\omterm{def:rc:counterparty-exposure:netting}{ensemble de compensation} et les trois CSA.
\item \textbf{Sens défavorable}~: \texttt{wrong\_way()}~; \texttt{fig\_rc\_exposure.py} produit les graphiques.
\end{tutsteps}
\textbf{Pour aller plus loin.} Faire payer à la banque le taux fixe sur le swap et observer comment le bénéfice de compensation
évolue~; construire le noyau en C++20 (\texttt{cpp/}) et le chronométrer par rapport à
l'agrégation Python sur 100\,000 trajectoires.
\end{tutorial}

\section{Construction~: le moteur d'exposition}\label{bld:rc:counterparty-exposure:exposure}

\begin{build}
\bfield{Rôle} Le moteur d'exposition de la société~: profils de chaque \omterm{def:rc:counterparty-exposure:netting}{ensemble de compensation} pour les limites de crédit
(PFE), pour les ajustements de valeur des chapitres~18 à~20 (EE, ENE), et pour le
moteur de risque du chapitre~29. \bfield{Interface}
\texttt{Market}, \texttt{simulate(market, horizon, steps\_per\_year,
paths, seed)} renvoyant \texttt{Scenarios} (\texttt{usd\_bond}, \texttt{eur\_bond})~;
transactions \texttt{Swap}, \texttt{XccySwap}, \texttt{FxForward} avec \texttt{values(sc)}~;
\texttt{collateral}, \texttt{exposure}, \texttt{profiles}, \texttt{epe},
\texttt{aggregate}, \texttt{conditional\_ee}. \texttt{firm::exposure::aggregate} en C++20
(\texttt{cpp/firm\_exposure.hpp}) et \texttt{firm\_exposure::aggregate} en Rust.
\bfield{Règles} Scénarios en probabilité neutre au risque~; un pas de la grille correspond à la période de marge en risque (dix
jours ouvrés)~; la dernière date de la grille, lorsque toutes les transactions sont arrivées à échéance, est exclue des
profils~; la PFE est le quantile inférieur à 97.5~\%.
\bfield{Tests d'acceptation} \texttt{code/firm/exposure/tests/}, \texttt{cpp/} et
\texttt{rust/}~: les trois implémentations reproduisent un même jeu de données à $10^{-9}$~; les transactions au pair
valent zéro~; les changes et les obligations actualisées sont des martingales~; la compensation et le collatéral réduisent le maximum de
l'EE, et une période de marge en risque plus longue l'augmente.
\bfield{Pour aller plus loin} Flux de trésorerie pendant la période de marge en risque~; taux EUR stochastiques et
corrélation taux--change~; réévaluation par régression des transactions remboursables par anticipation~; le noyau C++ appelé
depuis Python.
\end{build}

\begin{omsources}
\item Basel Committee on Banking Supervision, \emph{The standardised approach for measuring
counterparty credit risk exposures}, mars 2014.
\item Financial Crisis Inquiry Commission, \emph{The Financial Crisis Inquiry Report},
2011, notes du chapitre~20.
\item J.~Gregory, \emph{The xVA Challenge: Counterparty Risk, Funding, Collateral, Capital and Initial
Margin}, 4e édition, Wiley, 2020.
\end{omsources}

\section{Exercices}

\begin{exercise}[$\star$]\label{exo:rc:counterparty-exposure:1}
Pourquoi l'\omterm{def:rc:counterparty-exposure:profiles}{exposition attendue} d'un swap de taux d'intérêt est-elle en bosse, et celle d'un swap de devises
avec échange final est-elle croissante~?
\end{exercise}

\begin{exercise}[$\star$]\label{exo:rc:counterparty-exposure:2}
Deux transactions avec une même contrepartie valent $+10$ et $-7$ sur une trajectoire. Quelle est l'exposition
avec et sans \omterm{def:rc:counterparty-exposure:netting}{compensation avec déchéance du terme}~?
\end{exercise}

\begin{exercise}[$\star$]\label{exo:rc:counterparty-exposure:3}
Un CSA a un \omterm{def:rc:counterparty-exposure:csa}{seuil de collatéralisation} de 10~millions de dollars et un \omterm{def:rc:counterparty-exposure:csa}{montant minimal de transfert} (MTA) de 0.5~million de dollars~; l'\omterm{def:rc:counterparty-exposure:netting}{ensemble de compensation}
vaut 14.3~millions de dollars et la banque détient 4.0~millions de dollars. Calculer le montant que la banque appelle.
\end{exercise}

\begin{exercise}[$\star\star$]\label{exo:rc:counterparty-exposure:4}
Expliquer pourquoi doubler la période de marge en risque multiplie l'\omterm{def:rc:counterparty-exposure:profiles}{exposition attendue} collatéralisée par environ
$\sqrt2$.
\end{exercise}

\begin{exercise}[$\star\star$]\label{exo:rc:counterparty-exposure:5}
L'ENE compensée de l'\cref{ex:rc:counterparty-exposure:netting} est, en valeur absolue, plus grande que l'EE.
Expliquer ce que cela signifie pour la contrepartie et pour l'ajustement de valeur pour risque de crédit propre (DVA) de la banque (chapitre~18).
\end{exercise}

\begin{exercise}[$\star\star$]\label{exo:rc:counterparty-exposure:6}
Donner deux exemples de \omterm{def:rc:counterparty-exposure:wwr}{risque de corrélation défavorable} spécifique et un exemple de risque de corrélation favorable.
\end{exercise}

\begin{exercise}[$\star\star\star$]\label{exo:rc:counterparty-exposure:7}
\emph{Programmation.} Calculer l'EE maximale de l'\omterm{def:rc:counterparty-exposure:netting}{ensemble de compensation} collatéralisé avec une période de marge en risque de vingt jours et
vérifier le rapport à la valeur obtenue pour dix jours.
\end{exercise}

\begin{exercise}[$\star\star\star$]\label{exo:rc:counterparty-exposure:8}
\emph{Trouver la faille.} «~La contrepartie verse chaque jour une marge de variation avec un seuil nul,
donc notre exposition à son égard est nulle et elle n'a besoin d'aucune limite de crédit.~»
\end{exercise}

\section{Problème~: L'entreprise non collatéralisée}

\begin{problem}[{Problème du week-end --- signature d'un CSA}]\label{pb:rc:counterparty-exposure:1}
Une entreprise cliente paie le taux fixe à la banque sur un swap de taux de 50~millions de dollars à cinq ans (au taux au pair)
et achète à terme à la banque 20~millions d'euros dans un an (au taux à terme). Elle n'a pas de CSA~;
la banque lui en propose un avec un seuil nul et un \omterm{def:rc:counterparty-exposure:csa}{montant minimal de transfert} de 250\,000~dollars, ou
un autre avec un seuil de 5~millions de dollars.

\medskip\noindent\textbf{Partie I --- Sans CSA.}
\begin{enumerate}
\item Donner le taux au pair du swap et le strike du contrat à terme.
\item Donner la PFE maximale de l'\omterm{def:rc:counterparty-exposure:netting}{ensemble de compensation} et la date à laquelle elle est atteinte.
\item Donner son EPE sur un an et sur toute sa durée de vie.
\item Quelle transaction domine la première année, et pourquoi~?
\item Pourquoi une entreprise refuse-t-elle souvent de signer un CSA~?
\end{enumerate}

\medskip\noindent\textbf{Partie II --- Avec un CSA.}
\begin{enumerate}[resume]
\item Donner la PFE maximale et les EPE sous le CSA à seuil nul.
\item Pourquoi la PFE collatéralisée atteint-elle son maximum exactement à un an~?
\item Donner la PFE maximale et l'EPE sur la durée de vie avec le seuil de 5~millions de dollars.
\item Pourquoi le seuil de 5~millions de dollars ne change-t-il presque rien ici~?
\item Que faudrait-il à l'entreprise pour verser une marge de variation~?
\end{enumerate}

\medskip\noindent\textbf{Partie III --- Limites et fonds propres.}
\begin{enumerate}[resume]
\item La limite de la banque pour ce client est une PFE de 4~millions de dollars. Le client la respecte-t-il~?
\item Comment la limite modifie-t-elle la volonté de la banque de traiter davantage~?
\item Comment la banque réduirait-elle la PFE à la date de règlement-livraison~?
\item Que devient l'exposition si l'euro monte de 10~\% au cours du premier mois~?
\item Quelle partie est exposée à qui sur le contrat de change à terme si l'euro baisse~?
\end{enumerate}

\medskip\noindent\textbf{Partie IV --- Jugement.}
\begin{enumerate}[resume]
\item Pourquoi les banques continuent-elles de traiter sans collatéral avec les entreprises~?
\item Comment une période de marge en risque de vingt jours modifierait-elle les réponses~?
\item À quoi renonce la banque en acceptant un seuil~?
\item Énoncer le \emph{résultat de référence}~: la PFE maximale et l'EPE sur la durée de vie avant et après le
CSA à seuil nul.
\item En une phrase~: en quoi un CSA transforme-t-il une \omterm{def:rc:counterparty-exposure:exposure}{exposition de contrepartie}~?
\end{enumerate}
\end{problem}

\section{Questions d'entretien}

\begin{interviewq}[$\star$ \iqroles{risque, banque}]\label{iq:rc:counterparty-exposure:1}
Définir l'EE, la PFE et l'EPE. Laquelle sert à quoi~?
\end{interviewq}

\begin{interviewq}[$\star\star$ \iqroles{développeur, risque}]\label{iq:rc:counterparty-exposure:2}
Concevoir une simulation d'exposition pour l'ensemble du portefeuille de dérivés d'une banque. Où part le temps de calcul~?
\end{interviewq}

\begin{interviewq}[$\star\star$ \iqroles{risque}]\label{iq:rc:counterparty-exposure:3}
Qu'est-ce que la période de marge en risque, et pourquoi n'est-elle pas d'un jour pour un \omterm{def:rc:counterparty-exposure:netting}{ensemble de compensation}
avec appels de marge quotidiens~?
\end{interviewq}

\begin{interviewq}[$\star\star$ \iqroles{trader, risque}]\label{iq:rc:counterparty-exposure:4}
Donnez un exemple de \omterm{def:rc:counterparty-exposure:wwr}{risque de corrélation défavorable} et dites comment vous le mesureriez.
\end{interviewq}

\begin{interviewq}[$\star\star\star$ \iqroles{chercheur, développeur}]\label{iq:rc:counterparty-exposure:5}
Comment calculez-vous l'exposition d'une \omterm{def:rc:bermudans-and-callables:bermudan}{swaption bermudéenne} dans une simulation~?
\end{interviewq}

\begin{interviewq}[$\star\star\star$ \iqroles{risque, banque}]\label{iq:rc:counterparty-exposure:6}
Que sont devenues les contreparties de dérivés de Lehman dans les semaines qui ont suivi le 15 septembre 2008~?
\end{interviewq}
